\section{Image and text: MNIST-Sum details}
\label{app:image}

\subsection{Protocol}
\label{app:image:mnist}

\textbf{Corpus and representation.} Each example contains four native $28{\times}28$ MNIST digits
in raster order and the corresponding equation, including its sum. We sample $10^6$ examples and
hold out 10\% of the $10^4$ possible digit tuples. Pixels are binarised and flattened digit-major,
then row-major within each digit, so every 28-bit image position is one complete digit row. Text
belongs to a closed 40-symbol language (number words 0--36, ``$+$'', ``$=$'' and padding)
represented by deterministic 7-bit codes, four per 28-bit position. The sequence has 121
positions of 28 bits (3 text, 112 image and 6 markers; 3,388 bits in total). Neither modality uses a learned tokenizer.

\textbf{Training and evaluation.} The $12{\times}512$ denoiser (63.0M parameters) trains for 500K
steps with the shared diffusion objective, entropy controller, self-conditioning and task
mixture, and receives only the flattened sequence and its 1-D position. The positional settings
are those of the speech--text models, with four local Fourier features and an intra-patch period
of $P-1$; there is no segment embedding.
Evaluation uses EMA weights, 256 steps and $n{=}4{,}096$ independently sampled examples per
direction and split. Confidence intervals are 95\% Wilson intervals.

\subsection{Results, calibration and error decomposition}

\begin{table}[h]
\centering
\small
\setlength{\tabcolsep}{6pt}
\begin{tabular}{lrrr}
\toprule
& Seen tuples & Held-out tuples & CI (seen) \\
\midrule
Image$\rightarrow$text, exact string match & 97.56 & 97.85 & $[97.04, 97.99]$ \\
Text$\rightarrow$image, all four digits & 99.95 & 99.93 & $[99.82, 99.99]$ \\
Joint from noise, mutually consistent & 99.46 & --- & $[99.19, 99.65]$ \\
\bottomrule
\end{tabular}
\caption{MNIST-Sum task accuracy (\%), $n{=}4{,}096$ per direction and split. ``Held-out'' uses digit tuples
absent from training. Joint generation is counted correct only when the generated equation is
arithmetically valid and names the four digits in the generated image.}
\label{tab:app:mnist:complete}
\end{table}

\textbf{Image$\rightarrow$text} is scored by exact string equality, without a learned instrument.
The held-out interval $[97.36, 98.25]$ overlaps the seen interval almost entirely, so the
generalisation gap is not statistically distinguishable from zero. The seen result decomposes
into 97.66\% correct perception of all four digits, 99.90\% correct arithmetic given correct
perception and no malformed output. Per slot, perception is 99.41\%, 99.12\%, 99.58\% and
99.49\% for the top-left, top-right, bottom-left and bottom-right digits. An image-blind model
that always answers the most common sum scores 6.52\% on seen tuples and 7.80\% on held-out
tuples.

\textbf{Text$\rightarrow$image} and the image half of joint generation are read by one fixed
quadrant classifier, which is correct on 99.99\% of generated quadrants. Its ceiling depends on
which pool of real digits it is shown: it identifies all four digits in 96.09\% of composites
built from the MNIST test pool, from which the validation and held-out splits are drawn, and in
99.78\% of composites built from the train pool it was itself fitted on. Generated digits
imitate the train pool, so 99.78\% is the relevant ceiling; read against the 96.09\% figure the
model would appear to exceed the instrument, which is meaningless. These are reported as
calibrations rather than used to normalise the generation score.

\textbf{Joint co-generation} samples image and text from noise in one trajectory. Of the 4,096
pairs, 100.00\% have a well-formed equation, 99.71\% name the four digits that the classifier
reads in the generated image and 99.76\% are arithmetically valid; the 99.46\% that satisfy all
three is the reported rate. Cross-modal disagreement (0.29\%) and arithmetic error (0.24\%) are
of comparable size, so the residual failures cannot be attributed mainly to either.
